% !PRINTED 109
\weberpage{109}

Now $a + b\varrho$ is an integral or fractional number of the field $R(\varrho)$.
If we decompose the numerator and denominator of this number into their prime
factors, cancel common factors, and if we understand by $q_1, q_2, \ldots$ real
primes of the form $3N+2$, and by $\pi_1, \pi_1', \pi_2, \pi_2', \ldots$
conjugate pairs of imaginary primes of the form (6), then we obtain

\weq{8}{\begin{aligned}
a + b\varrho &= (-\varrho)^\lambda (\sqrt{-3})^n q_1^{s_1} q_2^{s_2} \ldots \pi_1^{t_1} \pi_1'^{t_1'} \pi_2^{t_2} \pi_2'^{t_2'} \ldots \\
a + b\varrho^2 &= (-\varrho^2)^\lambda (-\sqrt{-3})^n q_1^{s_1} q_2^{s_2} \ldots \pi_1^{t_1'} \pi_1'^{t_1} \pi_2^{t_2'} \pi_2'^{t_2} \ldots
\end{aligned}}

wherein

\[
\lambda, \ n, \ s_1, \ s_2 \ldots t_1, \ t_1', \ t_2, \ t_2' \ldots
\]

are positive or negative integers.

But now by (3) the product of the two numbers (8) is the cube of a rational
number, and this gives, since a rational number is decomposable into prime
factors in only one way, the following conditions:

\weq{8}{\begin{aligned}
n &\equiv 0, \quad s_1 \equiv 0, \quad s_2 \equiv 0 \ldots \\
t_1+t_1' &\equiv 0, \quad t_2+t_2' \equiv 0 \ldots \quad \text{(mod 3),}
\end{aligned}} % [sic: Weber repeats equation number (8); (9) follows.]

and we therefore set

\weq{9}{\begin{aligned}
n &= 3\mu, \quad s_1 = 3\sigma_1, \quad s_2 = 3\sigma_2 \ldots \\
t_1+t_1' &= 3\tau_1, \quad t_2+t_2' = 3\tau_2 \ldots,
\end{aligned}}

so that $n, \sigma_1, \sigma_2 \ldots \tau_1, \tau_2 \ldots$ are integers. % [sic: Weber prints n, although n=3\mu has just been set and \mu is evidently intended.]
Now by the formulae (7):

\weq{10}{\begin{aligned}
\pi^t \pi'^{t'} &= p^{-(t+t')} (\varrho,\eta)^{3t} (\varrho^2,\eta)^{3t'} \\
\pi^{t'} \pi'^{t} &= p^{-(t+t')} (\varrho,\eta)^{3t'} (\varrho^2,\eta)^{3t}.
\end{aligned}}

If finally we further denote by $\varepsilon$ a $9$th root of unity, then we can
set

\[
-\varrho = (-\varepsilon)^3
\]

and thereby represent $a+b\varrho$ and $a+b\varrho^2$ as cubes. We thus
obtain by (2):

\weq{11}{(\varrho, x_0)^3 = (-\varepsilon)^{3\lambda} (\sqrt{-3})^{3\mu} q_1^{3\sigma_1} q_2^{3\sigma_2} \ldots p_1^{-3\tau_1} p_2^{-3\tau_2} \ldots H_1^3,}

wherein

\weq{12}{H_1 = (\varrho,\eta_1)^{t_1} (\varrho^2,\eta_1)^{t_1'} (\varrho,\eta_2)^{t_2} (\varrho^2,\eta_2)^{t_2'} \ldots}

is a cyclotomic number. If we denote by $H_2$ the number that arises from
$H_1$ by the interchange of $\varrho$ with $\varrho^2$, then $H_1 H_2$ is a
rational number.

If we extract the third root from (11), then it follows, since a third root of
unity $\varrho_1$ may here appear as a factor besides,

\weq{13}{(\varrho, x_0) = (-\varepsilon)^\lambda \varrho_1 (\sqrt{-3})^\mu q_1^{\sigma_1} q_2^{\sigma_2} \ldots p_1^{-\tau_1} p_2^{-\tau_2} \ldots H_1,}
%CONT>>
